\subsection{Что показывают замеры}

\begin{itemize}
\item \textbf{Комбайнер} (наивный MapReduce $\to$ комбайнер, оба в одном воркере): время $20{,}87 \to 16{,}30$~с ($-22\%$), пик памяти $1\,234 \to 752$~МБ ($-39\%$). Без комбайнера маппер отдаёт мастеру каждое вхождение отдельной парой (ключ, 1) --- это $231\,216 + 1\,003\,894 + 3\,871\,492 = 5\,106\,602$ объектов, которые нужно сериализовать, передать между процессами и держать в памяти до свёртки. С комбайнером внутри сплита остаются только различные ключи.
\item \textbf{12 параллельных мапперов}: время $16{,}30 \to 10{,}96$~с, то есть ускорение всего в $1{,}5$ раза, а пик памяти вырос в $2{,}4$ раза ($1\,823$~МБ): 12 интерпретаторов Python, до 24 сплитов по 32~МиБ в работе и частичные счётчики. Ускорение далеко от 12 раз по закону Амдала: последовательными остаются чтение 26 чанков мастером, свёртка в 4 редьюсерах ($\approx1{,}5$~с), формирование и запись в S3 1,1~млн строк троек ($\approx3$~с) и запуск процессов на Windows. Сама map-фаза занимает лишь $\approx3{,}6$~с из 11.
\item \textbf{DuckDB против MapReduce}: $3{,}65$~с против $10{,}96$~с у лучшего MapReduce --- в 3 раза быстрее. Векторный движок на C++ параллельно разбирает CSV и строит хеш-агрегаты, без сериализации данных между процессами Python. Платит памятью ($1\,460$~МБ): параллельный разбор CSV держит в буферах крупные куски файла, а временные таблицы содержат 1,1~млн троек.
\item \textbf{Parquet вместо CSV} для DuckDB: время $3{,}65 \to 3{,}03$~с ($-17\%$), пик памяти $1\,460 \to 606$~МБ ($-58\%$). Parquet колоночный и в 5 раз меньше ($169$ против $848$~МБ): DuckDB читает 3 нужные колонки из 23 и не разбирает самое тяжёлое поле --- текст описания вакансии. Оставшиеся $\approx3$~с --- в основном постоянные затраты (запуск Python, загрузка httpfs, тройное соединение для троек).
\end{itemize}

\evidence{07_make_parquet.txt}{Создание Parquet-копии в hh-lake (make\_parquet.py)}

После профилирования результаты в \texttt{hh-lake} последний раз записали MapReduce на 12 воркерах и DuckDB по \emph{Parquet}; повторная сверка показывает, что и этот путь даёт те же числа:

\evidence{09_verify_after_profile.txt}{Сверка после профилирования: MapReduce (12 воркеров) против DuckDB по Parquet}

\begin{callout}
\textbf{Бонус: «до» и «после».}
\textbf{MapReduce}: $20{,}87$~с / $1\,234$~МБ $\to$ $10{,}96$~с / $1\,823$~МБ (в $1{,}9$ раза быстрее, ценой памяти) или $16{,}30$~с / $752$~МБ (в $1{,}3$ раза быстрее и на 39\% меньше памяти, если ресурс --- память).
\textbf{DuckDB}: $3{,}65$~с / $1\,460$~МБ $\to$ $3{,}03$~с / $606$~МБ (на 17\% быстрее и в $2{,}4$ раза меньше памяти).
Результаты всех вариантов совпадают до md5.
\end{callout}
